% !TeX encoding = UTF-8
% !TeX spellcheck = sk_SK
\documentclass[]{tukediphc}
%% -----------------------------------------------------------------
%% tento subor ma kodovanie utf-8
%%
%% na kompilaciu pouzivajte format pdflatex 
%%
%% V pripade problemov kontaktujte Jána Bušu st. (jan.busa@tuke.sk)
%%
%% November 2015
%% -----------------------------------------------------------------
%%
%\usepackage[dvips]{graphicx}
%\DeclareGraphicsExtensions{.eps}
\usepackage[pdftex]{graphicx}
\DeclareGraphicsExtensions{.pdf,.png,.jpg,.mps}
\graphicspath{{figures/}} % priecinok na obrazky
%%


%\usepackage[utf8]{inputenc}  % je v cls-subore
%\usepackage[T1]{fontenc}  % je v cls-subore
\usepackage{lmodern,textcase}
\usepackage[slovak]{babel}
\def\refname{Zoznam použitej literatúry}
\usepackage{latexsym}
\usepackage{dcolumn} % zarovnanie cisiel v tabulke podla des. ciarky
\usepackage{hhline}
\usepackage{amsmath,amsfonts,amssymb}
\usepackage{nicefrac} % pekne zlomky
\IfFileExists{upgreek.sty}{%
  \usepackage{upgreek}% napr. $\upmu\mathrm{m}$ pre mikrometer ...
}{%
  \newcommand{\upmu}{\mu}%
}
\usepackage[final]{showkeys}%color%notref%notcite%final
\usepackage[slovak,noprefix]{nomencl}
\makeglossary % prikaz na vytvorenie suboru .glo

% Bezpecnost: ak sa niektore makra z triedy nenacitaju, nech to nepadne.
\providecommand{\abstraktsk}{}
\providecommand{\abstrakteng}{}
\providecommand{\kabstrakt}{}
\providecommand{\zadanieprace}{}
\providecommand{\cestnevyhlasenie}{}
\providecommand{\podakovanie}{}
\providecommand{\kpodakovania}{}
\providecommand{\curriculumvitae}{}
\providecommand{\kcurriculumvitae}{}


% Pouzit v pripade velkeho poctu subsection v tableofcontents
%\makeatletter
%\renewcommand*\l@subsection{\@dottedtocline{2}{1.5em}{3.5em}}
%\newcommand*\l@subsection{\@dottedtocline{2}{1.5em}{2.3em}}
%\newcommand*\l@subsubsection{\@dottedtocline{3}{3.8em}{3.2em}}
%\makeatother


%\def\thefigure{\Roman{section}.\arabic{figure}}

%\usepackage{parskip}% 'zhusti' polozky obsahu
%% Cislovane citovanie
\usepackage[numbers]{natbib}
%%
%% Citovanie podľa mena autora a roku
%\usepackage{natbib} \citestyle{chicago}
% -----------------------------------------------------------------
%% tlač !!!
\usepackage[pdftex,unicode=true,bookmarksnumbered=true,
bookmarksopen=true,pdfmenubar=true,pdfview=Fit,linktocpage=true,
pageanchor=true,bookmarkstype=toc,pdfpagemode=UseOutlines,
pdfstartpage=1]{hyperref}
\hypersetup{%
baseurl={http://www.tuke.sk/sevcovic},
pdfcreator={pdfcsLaTeX},
pdfkeywords={dochádzka, príchody, odchody, zamestnanci, bezkontaktné},
pdftitle={15. Systém na zaznamenávanie príchodov a odchodov zamestnancov bezkontaktne},
pdfauthor={Martin Oros},
pdfsubject={Zadanie ZSI}
} 
%% nehodiace zakomentujte !
\dippraca{Zadanie ZSI}
%\bakpraca{Bakalárska práca}
%%
\nazov{15. Systém na zaznamenávanie príchodov a odchodov zamestnancov bezkontaktne}
%% ked praca nema 'podnazov' zakomentujte nasledujuci riadok
%% alebo polozku nechajte prazdnu
\podnazov{}
\jazyk{Slovenský}
% anglicky nazov
\title{Contactless employee arrivals and departures logging system}
\autor{Martin Oros}
\veduciprace{doc.~Ing.~Vojtech~Čierny, CSc.}
%\konzultanta{Ing.~Matej~Biely, PhD.}
%\konzultantb{RNDr.~Marián~Čierny, DrSc.}
\titul{Bc}
\univerzita{Technická univerzita v~Košiciach}
\fakulta{Fakulta elektrotechniky a informatiky}
\skratkafakulty{FEI}
\katedra{Katedra umelej inteligencie (KUI)}
\skratkakatedry{KUI}
\odbor{}
\specializacia{}
\abstrakt{Abstrakt je povinnou súčasťou každej práce. Je výstižnou
charakteristikou obsahu dokumentu. Nevyjadruje hodnotiace stanovisko
autora. Má byť taký informatívny, ako to povoľuje podstata práce.
Text abstraktu sa píše ako jeden odstavec. Abstrakt neobsahuje odkazy
na samotný text práce. Mal by mať rozsah 250 až 500 slov. Pri
štylizácii sa používajú celé vety, slovesá v činnom rode a tretej
osobe. Používa sa odborná terminológia, menej zvyčajné termíny,
skratky a~symboly sa pri prvom výskyte v texte definujú.}
\klucoveslova{dochádzka, príchody, odchody, zamestnanci, bezkontaktné}
\abstrakte{Text abstraktu v~svetovom jazyku je potrebný pre integráciu
do medzinárodných informačných systémov. Ak nie je možné cudzojazyčnú
verziu abstraktu umiestniť na jednej strane so slovenským abstraktom,
je potrebné umiestniť ju na samostatnú stranu (cudzojazyčný abstrakt
nemožno deliť a~uvádzať na dvoch strabách).}
\keywords{attendance, arrivals, departures, employees, contactless}
\datumodovzdania{11.~4.~2026}
\datumobhajoby{15.~6.~2026}
\mesto{Košice}
\pocetstran{\pageref{page:posledna}}
\kategoria{Záverečná práca}

\begin{document}
\renewcommand{\figurename}{Obrázok}	
\renewcommand\theHfigure{\theHsection.\arabic{figure}}
\renewcommand\theHtable{\theHsection.\arabic{table}}
\bibliographystyle{dcu}

\prvastrana

%\titulnastrana

%\analytickylist


%\errata % zaciatok erraty
%Ak je potrebné, autor na tomto mieste opraví chyby, ktoré našiel po
%vytlačení práce. Opravy sa uvádzajú takým písmom, akým je napísaná
%práca. Ak zistíme chyby až po vytlačení a zviazaní práce, napíšem
%erráta na samostatný lístok, ktorý vložíme na toto miesto. Najlepšie je
%lístok prilepiť \citep{kat}.
%
%Forma:
%
%%\tabcolsep=10pt
%\begin{table}[!hb]
%	\centering
%	\begin{tabular}{|c|c|c|c|}\hline
%Strana & Riadok & Chybne & Správne \\\hline\hline
%12 & 6 & publikácia & prezentácia \\\hline
%22 & 23 & internet & intranet \\\hline
%& & & \\\hline
%& & & \\\hline
%	\end{tabular}
%\end{table}
%\kerrata % koniec erraty

%\abstraktsk % abstrakt v SK 

\abstrakteng % abstrakt v ENG

\kabstrakt % koniec abstraktov, nova strana

% Na tomto mieste bude vložené zadanie diplomovej práce
\zadanieprace

\cestnevyhlasenie
% Niektorí autori metodických príručiek o~záverečných prácach sa
% nazdávajú, že takéto vyhlásenie je zbytočné, nakoľko povinnosť
% vypracovať záverečnú prácu samostatne, vyplýva študentovi zo zákona a
% na autora práce sa vzťahuje autorský zákon.

% Poďakovanie (odstránené na požiadanie)
%\podakovanie
%Na tomto mieste môže byť vyjadrenie poďakovania napr. vedúcemu
%diplomovej práce, resp. konzultantom, za pripomienky a~odbornú pomoc
%pri vypracovaní diplomovej práce.
%
%Na tomto mieste môže byť vyjadrenie poďakovania napr. vedúcemu
%diplomovej práce, respektíve konzultantom, za pripomienky a~odbornú
%pomoc pri vypracovaní diplomovej práce.
%
%Na tomto mieste môže byť vyjadrenie poďakovania napr. vedúcemu
%diplomovej práce alebo konzultantom za pripomienky a~odbornú pomoc pri
%vypracovaní diplomovej práce.
%\kpodakovania

% Predhovor (odstránené na požiadanie)
%\predhovor
%Predhovor je povinnou náležitosťou záverečnej práce, pozri
%\citep{gonda}. V~predhovore autor uvedie základné charakteristiky
%svojej záverečnej práce a~okolnosti jej vzniku. Vysvetlí dôvody, ktoré
%ho viedli k~voľbe témy, cieľ a~účel práce a~stručne informuje
%o~hlavných metódach, ktoré pri spracovaní záverečnej práce použil.
%\kpredhovoru

\thispagestyle{empty}
\tableofcontents
\newpage

% Zoznam obrázkov (odstránené na požiadanie)
%\thispagestyle{empty}
%{	\makeatletter
%	\renewcommand{\l@figure}{\@dottedtocline{1}{1.5em}{3.5em}}
%	\makeatother
%	\listoffigures}
%\newpage

\thispagestyle{empty}
%\addcontentsline{toc}{section}{\numberline{}Zoznam tabuliek}
\listoftables
\newpage

\thispagestyle{empty}
%\addcontentsline{toc}{section}{\numberline{}Zoznam symbolov a
%skratiek}
\printglossary % vlozenie zoznamu skratiek a symbolov
\newpage

% Slovník termínov (odstránené na požiadanie)
%\slovnikterminov
%\begin{description}
%	\item[Dizertácia] ...
%\end{description}
%\kslovnikterminov
%
% Úvod (odstránené na požiadanie)
%\include{uvod}
%
% !TeX encoding = UTF-8
% !TeX spellcheck = sk_SK
% !TeX root=tukedip.tex
\section{Formulácia úlohy}
Sledovanie 
dochádzky 
zamestnancov - systém na zaznamenávanie príchodov a odchodov zamestnancov bezkontaktne.
%
% !TeX root=tukedip.tex
% !TeX encoding = UTF-8
% !TeX spellcheck = sk_SK
%
% Required packages (add to preamble if not already present):
%   \usepackage{booktabs}
%   \usepackage{tabularx}
%   \usepackage{plantuml}
%   \usepackage{lscape}
%   \usepackage{array}
%

\section{Analýza požiadaviek (Requirements Analysis)}

Táto kapitola sa venuje analýze požiadaviek na webovú aplikáciu pre bezkontaktné
zaznamenávanie príchodov a~odchodov zamestnancov pomocou QR kódov. Systém pozostáva
z~dvoch hlavných rozhraní: rozhrania pre zamestnancov, ktoré umožňuje prihlásenie
a~skenovanie QR kódu na zaznamenanie príchodu alebo odchodu, a~administrátorského
rozhrania poskytujúceho správu záznamov, export dát, filtrovanie histórie dochádzky
a~generovanie QR kódov pre zamestnancov. Analýza je štruktúrovaná prostredníctvom
používateľských príbehov a~prípadov použitia, ktoré tvoria základ pre návrh
a~implementáciu systému.

% -----------------------------------------------------------------------
\subsection{Používateľské príbehy (User Stories)}

Používateľské príbehy opisujú požadované funkcie systému z~pohľadu jednotlivých
aktérov -- zamestnanca a~administrátora. Každý príbeh obsahuje akceptačné kritérium,
ktoré definuje podmienky úspešného splnenia danej požiadavky. Prehľad príbehov je
uvedený v~Tabuľke~\ref{t:user-stories}.

\renewcommand{\arraystretch}{1.12}
\setlength{\tabcolsep}{3pt}
\footnotesize
\begin{longtable}{|p{1.2cm}|p{2.6cm}|p{5.4cm}|p{5.4cm}|}
\caption{Používateľské príbehy}\label{t:user-stories}\\
\hline
\textbf{ID} & \textbf{Rola} & \textbf{Príbeh} & \textbf{Akceptačné kritérium} \\
\hline
\endfirsthead
\hline
\textbf{ID} & \textbf{Rola} & \textbf{Príbeh} & \textbf{Akceptačné kritérium} \\
\hline
\endhead
US-01 & Zamestnanec &
Ako zamestnanec chcem sa prihlásiť do systému pomocou mena a~hesla, aby som mal
prístup k~funkcii skenovania QR kódu. &
Prihlásenie prebehne úspešne len pri zadaní správnych prihlasovacích údajov;
nesprávne údaje zobrazia chybovú správu. \\
\hline
US-02 & Zamestnanec &
Ako zamestnanec chcem naskenovať svoj QR kód pomocou kamery zariadenia, aby bol
môj príchod alebo odchod automaticky zaznamenaný. &
Systém rozpozná QR kód, určí správny typ záznamu (príchod/odchod) a~uloží ho
s~časovou pečiatkou. \\
\hline
US-03 & Zamestnanec &
Ako zamestnanec chcem po úspešnom skene vidieť potvrdenie o~zaznamenaní príchodu
alebo odchodu, aby som bol istý, že dáta boli uložené. &
Na obrazovke sa zobrazí správa s~typom záznamu a~aktuálnym časom do 2~sekúnd
od skenovania. \\
\hline
US-04 & Zamestnanec &
Ako zamestnanec chcem sa odhlásiť zo systému, aby bola moja relácia bezpečne
ukončená. &
Po odhlásení je relácia zrušená a~používateľ je presmerovaný na prihlasovaciu
stránku. \\
\hline
US-05 & Administrátor &
Ako administrátor chcem sa prihlásiť do oddeleného administrátorského rozhrania
pomocou svojich prihlasovacích údajov, aby som mal prístup k~správe systému. &
Administrátorský účet umožňuje prístup iba do admin dashboardu; bežní zamestnanci
nemôžu pristupovať k~admin rozhraniu. \\
\hline
US-06 & Administrátor &
Ako administrátor chcem prehliadať históriu dochádzky všetkých zamestnancov
s~možnosťou filtrovania podľa mena, dátumu a~typu záznamu, aby som prehľadne videl
dochádzku. &
Záznamy sú zobrazené v~tabuľke; filtrovanie vracia iba relevantné záznamy
zodpovedajúce zvoleným kritériám. \\
\hline
US-07 & Administrátor &
Ako administrátor chcem upravovať alebo mazať záznamy dochádzky, aby som mohol
opraviť prípadné chyby v~dátach. &
Zmenené záznamy sú okamžite aktualizované v~databáze; zmazané záznamy sú trvalo
odstránené po potvrdení. \\
\hline
US-08 & Administrátor &
Ako administrátor chcem exportovať záznamy dochádzky do súboru Excel, aby som
mohol dáta ďalej spracovávať alebo archivovať mimo systému. &
Exportovaný súbor obsahuje všetky zobrazené záznamy (prípadne filtrované) vo
formáte XLSX a~je automaticky stiahnutý do zariadenia. \\
\hline
US-09 & Administrátor &
Ako administrátor chcem generovať QR kódy pre jednotlivých zamestnancov, aby každý
zamestnanec mal unikátny kód na skenovanie. &
Vygenerovaný QR kód je jednoznačne priradený k~danému zamestnancovi a~je možné
ho stiahnuť alebo vytlačiť. \\
\hline
US-10 & Administrátor &
Ako administrátor chcem spravovať zoznam zamestnancov (pridávanie, deaktivácia),
aby systém odrážal aktuálny stav personálu. &
Nový zamestnanec je po pridaní okamžite dostupný v~systéme; deaktivovaný
zamestnanec nemôže skenovať QR kód. \\
\hline
\end{longtable}

% -----------------------------------------------------------------------
\subsection{Prípady použitia (Use Cases)}

Prípady použitia podrobnejšie opisujú interakcie aktérov so systémom vrátane
predpodmienok a~očakávaných výsledkov. Prehľad prípadov použitia je uvedený
v~Tabuľke~\ref{t:use-cases}.

\renewcommand{\arraystretch}{1.12}
\setlength{\tabcolsep}{3pt}
\footnotesize
\begin{longtable}{|p{1.1cm}|p{3.2cm}|p{6.2cm}|p{4.2cm}|}
\caption{Prípady použitia}\label{t:use-cases}\\
\hline
\textbf{ID} & \textbf{Názov} & \textbf{Popis} & \textbf{Poznámka} \\
\hline
\endfirsthead
\hline
\textbf{ID} & \textbf{Názov} & \textbf{Popis} & \textbf{Poznámka} \\
\hline
\endhead
UC-01 & Prihlásenie zamestnanca &
Zamestnanec zadá svoje meno a~heslo na prihlasovacej stránke aplikácie. &
Aktér: Zamestnanec. Pred: platný účet. Výs: prihlásený, presmerovaný na skenovanie. \\
\hline
UC-02 & Skenovanie QR kódu &
Zamestnanec nasmeruje kameru zariadenia na vytlačený QR kód; systém ho
automaticky rozpozná pomocou knižnice jsQR.js. &
Aktér: Zamestnanec. Pred: prihlásený, funkčná kamera. Výs: uložený príchod/odchod s~časovou pečiatkou. \\
\hline
UC-03 & Zobrazenie potvrdenia &
Po úspešnom skene systém zobrazí potvrdzovaciu správu s~typom záznamu a~časom. &
Aktér: Zamestnanec. Pred: úspešný sken (UC-02). Výs: vizuálne potvrdenie záznamu. \\
\hline
UC-04 & Odhlásenie &
Používateľ klikne na tlačidlo odhlásenia; systém zruší reláciu. &
Aktér: Zamestnanec/Administrátor. Pred: prihlásený. Výs: relácia zrušená, presmerovaný na prihlásenie. \\
\hline
UC-05 & Prihlásenie administrátora &
Administrátor zadá svoje prihlasovacie údaje na osobitnej admin prihlasovacej
stránke. &
Aktér: Administrátor. Pred: platný admin účet. Výs: prihlásený, presmerovaný na admin dashboard. \\
\hline
UC-06 & Prehliadanie a~filtrovanie dochádzky &
Administrátor prezerá tabuľku dochádzky a~aplikuje filtre podľa mena zamestnanca,
dátumového rozsahu alebo typu záznamu. &
Aktér: Administrátor. Pred: prihlásený (UC-05). Výs: zobrazené záznamy podľa filtra. \\
\hline
UC-07 & Úprava záznamu dochádzky &
Administrátor vyberie konkrétny záznam a~upraví jeho hodnoty (typ, čas) cez
editačný formulár. &
Aktér: Administrátor. Pred: prihlásený, záznam existuje. Výs: údaje uložené a~zobrazené. \\
\hline
UC-08 & Vymazanie záznamu dochádzky &
Administrátor vyberie záznam na vymazanie a~potvrdí akciu. &
Aktér: Administrátor. Pred: prihlásený, záznam existuje. Výs: záznam odstránený. \\
\hline
UC-09 & Export do Excelu &
Administrátor klikne na tlačidlo exportu; systém vygeneruje XLSX súbor
s~aktuálne zobrazenými záznamami. &
Aktér: Administrátor. Pred: prihlásený, existujú záznamy. Výs: XLSX stiahnutý. \\
\hline
UC-10 & Generovanie QR kódu &
Administrátor vyberie zamestnanca a~nechá systém vygenerovať jeho unikátny QR
kód pomocou knižnice chillerlan/php-qrcode. &
Aktér: Administrátor. Pred: prihlásený, zamestnanec evidovaný. Výs: QR kód vygenerovaný a dostupný. \\
\hline
UC-11 & Správa zamestnancov &
Administrátor spravuje zoznam zamestnancov (pridanie, úprava, deaktivácia), aby systém
odrážal aktuálny stav personálu. &
Aktér: Administrátor. Pred: prihlásený. Výs: zmeny uložené; deaktivovaný zamestnanec nemôže skenovať. \\
\hline
\end{longtable}

% -----------------------------------------------------------------------
\subsection{Use Case Diagram}

Diagram prípadov použitia (Obrázok~\ref{o:use-case-diagram}) znázorňuje vzťahy medzi
aktérmi systému a~jednotlivými funkciami aplikácie. Systém definuje dvoch aktérov:
\textit{Zamestnanca}, ktorý využíva funkcie prihlasovania a~skenovania QR kódu, a~\textit{Administrátora}, ktorý spravuje záznamy dochádzky, generuje QR kódy a~exportuje
dáta. Prípad použitia \textit{Skenovanie QR kódu} obsahuje \texttt{<<include>>} vzťah
na \textit{Zobrazenie potvrdenia}, keďže potvrdenie je neoddeliteľnou súčasťou
procesu skenovania. Prípad použitia \textit{Prehliadanie dochádzky} je rozšírený
prostredníctvom \texttt{<<extend>>} o~\textit{Filtrovanie záznamov}, ktoré je
voliteľnou, ale bežne využívanou funkciou. Podobne \textit{Export do Excelu} rozširuje
\textit{Prehliadanie dochádzky} ako voliteľná akcia nadväzujúca na zobrazenie záznamov.

\begin{figure}[!ht]
\centering
\begin{plantuml}
@startuml
left to right direction
skinparam packageStyle rectangle
skinparam actorStyle awesome
skinparam usecase {
  BackgroundColor #EEF4FB
  BorderColor #2C6DAD
  ArrowColor #2C6DAD
  ActorBorderColor #1A3F6F
  ActorBackgroundColor #D6E8FA
}

actor "Zamestnanec" as EMP
actor "Administrátor" as ADM

rectangle "QR Attendance Tracker" {

  usecase "Prihlásenie\nzamestnanca" as UC1
  usecase "Skenovanie\nQR kódu" as UC2
  usecase "Zobrazenie\npotvrdenia" as UC3
  usecase "Odhlásenie" as UC4

  usecase "Prihlásenie\nadministrátora" as UC5
  usecase "Prehliadanie\ndochádzky" as UC6
  usecase "Filtrovanie\nzáznamov" as UC7
  usecase "Úprava záznamu" as UC8
  usecase "Vymazanie záznamu" as UC9
  usecase "Export do Excelu" as UC10
  usecase "Generovanie\nQR kódu" as UC11

}

EMP --> UC1
EMP --> UC2
EMP --> UC4

UC2 .> UC3 : <<include>>

ADM --> UC5
ADM --> UC6
ADM --> UC8
ADM --> UC9
ADM --> UC10
ADM --> UC11
ADM --> UC4

UC7 .> UC6 : <<extend>>
UC10 .> UC6 : <<extend>>

@enduml
\end{plantuml}
\caption{Diagram prípadov použitia systému}\label{o:use-case-diagram}
\end{figure}

% -----------------------------------------------------------------------
\subsection{Prioritizácia požiadaviek (MoSCoW)}

V tejto časti sú požiadavky z~pohľadu používateľských príbehov a~prípadov použitia
zoraďované pomocou metódy MoSCoW do kategórií:
\textbf{M} (Must have), \textbf{S} (Should have), \textbf{C} (Could have),
\textbf{W} (Won't have).

\renewcommand{\arraystretch}{1.12}
\setlength{\tabcolsep}{3pt}
\footnotesize
\begin{longtable}{|p{1.3cm}|p{1.3cm}|p{9.0cm}|p{1.0cm}|}
\caption{MoSCoW prioritizácia požiadaviek}\label{t:moscow}\\
\hline
\textbf{US} & \textbf{UC} & \textbf{Požiadavka (stručne)} & \textbf{Prio} \\
\hline
\endfirsthead
\hline
\textbf{US} & \textbf{UC} & \textbf{Požiadavka (stručne)} & \textbf{Prio} \\
\hline
\endhead
US-01 & UC-01 & Prihlásenie zamestnanca (meno/heslo) pre prístup ku skenovaniu. & M \\
\hline
US-02 & UC-02 & Skenovanie QR kódu a automatické uloženie príchodu/odchodu s časom. & M \\
\hline
US-03 & UC-03 & Zobrazenie potvrdenia po úspešnom skene. & S \\
\hline
US-04 & UC-04 & Odhlásenie používateľa a ukončenie relácie. & M \\
\hline
US-05 & UC-05 & Prihlásenie administrátora do admin rozhrania. & M \\
\hline
US-06 & UC-06 & Prehliadanie histórie dochádzky a filtrovanie záznamov. & M \\
\hline
US-07 & UC-07 & Úprava záznamu dochádzky (typ/čas). & S \\
\hline
US-07 & UC-08 & Vymazanie záznamu dochádzky (s potvrdením). & S \\
\hline
US-08 & UC-09 & Export záznamov do XLSX. & C \\
\hline
US-09 & UC-10 & Generovanie QR kódov pre zamestnancov. & M \\
\hline
US-10 & UC-11 & Správa zamestnancov (pridanie/deaktivácia). & S \\
\hline
\end{longtable}

%
\include{jadroprace}
%
\include{zaver}
%
\include{literatura}
%
\include{prilohy}
%
\include{prilohaa}
%
\include{prilohab}
%
\include{prilohac}
%
% zivotopis autora
\newpage
\phantomsection
\protect\label{page:posledna}
\curriculumvitae
Táto časť je nepovinná. Autor tu môže uviesť svoje biografické
údaje, údaje o~záujmoch, účasti na~projektoch, účasti na~súťažiach,
získané ocenenia, zahraničné pobyty na~praxi, domácu prax, publikácie
a~pod.
\kcurriculumvitae

\end{document}
%%